\chapter*{Introduction}
\addcontentsline{toc}{chapter}{Introduction}

\section*{Background}
Air quality is one of the most important factors in human quality of life. Ukraine has historically faced air quality challenges attributable to two principal sources: emissions from heavy industrial facilities — which account for over 60\% of total national emissions \cite{pehchevski2020ukraine} — and uncontrolled agricultural burning, including the burning of crop residues. Since the onset of the full-scale military conflict in February 2022, a third significant source has emerged: fires caused by missile and drone strikes. Research covering the first two years of the conflict demonstrates that fire-related emissions became a major contributor to air pollution\cite{MALYTSKA2026108959}.

Numerous resources study air quality using a variety of approaches — from ground-level sensor networks to satellite technologies. The resulting data is used both in scientific research and in routine public notification systems. While deploying large-scale sensor networks was once technically demanding, advances in embedded systems and network infrastructure have made continuous, high-frequency monitoring across hundreds of stations increasingly accessible — generating data volumes and velocities that conventional processing approaches are ill-suited to handle within the narrow time window in which pollution data retains its actionable value, given that acute spikes in PM2.5 or NO\textsubscript{2} typically persist for hours rather than days.

\section*{Problem Statement}
While applications such as Kyiv Digital and EcoCity include built-in notification systems for air quality degradation across various cities, none of the systems examined provides consistent coverage across all regions and at least district-level population centers, based on the publicly available features reviewed at the time of writing. Furthermore, all of these systems operate on the basis of hourly data readings, meaning that notifications may be delayed — an outcome that can affect public health, particularly for individuals with respiratory conditions or other vulnerabilities. Coverage and timeliness are the principal guarantees of effective early warning against negative health impacts.

\section*{Research Objectives}
This work aims to address these gaps by designing and implementing a cloud-native streaming data pipeline for air quality monitoring across 414 Ukrainian cities, capable of processing high-throughput sensor readings with sub-minute end-to-end latency, surfacing real-time alerts and dashboard visualizations, and persisting both raw sensor readings and aggregated results for historical analysis.

\section*{Scope and Limitations}
This research focuses specifically on data ingestion, processing, alerting and visualization infrastructure to support real-time data analytics. The implementation prioritizes cloud-native solutions, leveraging modern cloud platforms that offer elastic scalability, managed services, reduced operational overhead, and pay-as-you-go cost models that have become industry standards for data-intensive applications. The study deliberately excludes application business logic development, detailed examination of traditional relational database implementations, and batch processing methodologies, as these fall outside the core streaming architecture objectives.

\section*{Methodology}
This research follows a systematic approach beginning with high-level system architecture design encompassing all components from data-generating applications through visualization layers. Subsequently, we evaluate and select implementation tools, emphasizing cloud-managed services across storage, processing, and analytics domains. The implementation phase involves system construction, iterative debugging, and comprehensive performance testing. Finally, we analyze achieved outcomes, throughput metrics, latency measurements, and implementation challenges encountered during development.

\section*{Significance of the Study}
This work contributes a practical reference implementation of a streaming data pipeline suitable for integration into a large-scale public air quality notification system. The architectural patterns and implementation details documented here provide a concrete foundation for building real-time data processing and analytics systems as monitoring infrastructure continues to develop. The work explicitly addresses the limitations of existing implementations — specifically their limited geographic coverage and hourly notification latency — and demonstrates how these gaps can be closed through modern streaming architecture.